\documentclass[12pt, a4paper]{article}

\usepackage[margin=1in]{geometry}
\usepackage{amsmath}
\usepackage{booktabs}
\usepackage{hyperref}
\usepackage{listings}
\usepackage{xcolor}
\usepackage{graphicx}
\usepackage{enumitem}
\usepackage{fancyhdr}
\usepackage{titlesec}
\usepackage{array}

\hypersetup{
    colorlinks=true,
    linkcolor=blue,
    urlcolor=blue,
    citecolor=blue
}

\lstset{
    basicstyle=\ttfamily\small,
    backgroundcolor=\color{gray!10},
    frame=single,
    breaklines=true,
    keywordstyle=\color{blue},
    commentstyle=\color{green!60!black},
    stringstyle=\color{red!70!black},
    language=Python,
    showstringspaces=false,
    numbers=left,
    numberstyle=\tiny\color{gray},
    numbersep=5pt,
    tabsize=4
}

\pagestyle{fancy}
\fancyhf{}
\rhead{Fraud Detection Project}
\lhead{\leftmark}
\rfoot{Page \thepage}

\title{
    \vspace{-1cm}
    \textbf{Financial Transaction Fraud Detection} \\[0.5em]
    \large Project Report \\[0.3em]
    \normalsize Using Logistic Regression and Streamlit
}
\author{AIML Project}
\date{\today}

\begin{document}

\maketitle
\thispagestyle{fancy}
\tableofcontents
\newpage

% ─────────────────────────────────────────────
\section{Project Overview}
% ─────────────────────────────────────────────

This project builds an end-to-end machine learning pipeline to detect fraudulent financial transactions. The workflow consists of three stages:

\begin{enumerate}[label=\textbf{\arabic*.}]
    \item \textbf{Exploratory Data Analysis (EDA)} — understanding class distribution, transaction patterns, and suspicious behaviour signals.
    \item \textbf{Feature Engineering and Model Training} — constructing derived features and training a Logistic Regression classifier using a scikit-learn \texttt{Pipeline}.
    \item \textbf{Deployment} — serving predictions through an interactive Streamlit web application.
\end{enumerate}

% ─────────────────────────────────────────────
\section{Dataset Description}
% ─────────────────────────────────────────────

\subsection{Source}

The dataset (\texttt{AIML Dataset.csv}) is a synthetic mobile money transaction log containing \textbf{6,362,620 rows} and \textbf{11 columns}, with no missing values.

\subsection{Columns}

\begin{center}
\begin{tabular}{lll}
\toprule
\textbf{Column} & \textbf{Type} & \textbf{Description} \\
\midrule
\texttt{step}           & int64   & Time unit (1 step = 1 hour) \\
\texttt{type}           & str     & Transaction type (PAYMENT, TRANSFER, etc.) \\
\texttt{amount}         & float64 & Transaction amount \\
\texttt{nameOrig}       & str     & Sender account ID \\
\texttt{oldbalanceOrg}  & float64 & Sender balance before transaction \\
\texttt{newbalanceOrig} & float64 & Sender balance after transaction \\
\texttt{nameDest}       & str     & Receiver account ID \\
\texttt{oldbalanceDest} & float64 & Receiver balance before transaction \\
\texttt{newbalanceDest} & float64 & Receiver balance after transaction \\
\texttt{isFraud}        & int64   & Target label (1 = fraud, 0 = legitimate) \\
\texttt{isFlaggedFraud} & int64   & System-flagged (not used in model) \\
\bottomrule
\end{tabular}
\end{center}

\subsection{Class Imbalance}

Only \textbf{0.13\%} of all transactions are fraudulent — a severe class imbalance that is a key challenge for the classifier.

\begin{align}
\text{Fraud rate} = \frac{|\{i : \texttt{isFraud}_i = 1\}|}{N} \times 100 = 0.13\%
\end{align}

% ─────────────────────────────────────────────
\section{Exploratory Data Analysis}
% ─────────────────────────────────────────────

\subsection{Transaction Type Distribution}

The dataset contains five transaction types. A bar chart of their counts reveals that \textbf{CASH\_OUT} and \textbf{PAYMENT} are the most frequent types.

\subsection{Fraud by Transaction Type}

A critical finding is that \textbf{fraud occurs exclusively in TRANSFER and CASH\_OUT} transactions. No fraudulent activity was observed in PAYMENT, DEBIT, or CASH\_IN.

\begin{center}
\begin{tabular}{lrr}
\toprule
\textbf{Type} & \textbf{Total Count} & \textbf{Has Fraud?} \\
\midrule
CASH\_OUT  & 2,237,500 & Yes \\
TRANSFER   &   532,909 & Yes \\
PAYMENT    & (remaining) & No \\
DEBIT      & (remaining) & No \\
CASH\_IN   & (remaining) & No \\
\bottomrule
\end{tabular}
\end{center}

\subsection{Amount Distribution}

Transaction amounts are heavily right-skewed. A log-transform (\texttt{log1p}) is applied for visualisation:

\begin{lstlisting}
sns.histplot(np.log1p(df["amount"]), bins=100, kde=True, color="green")
\end{lstlisting}

\subsection{Zero-Balance-After-Transfer Pattern}

A suspicious pattern was identified: transactions where the sender had a positive balance before the transaction but \textit{zero} balance after, specifically in TRANSFER or CASH\_OUT. This accounts for \textbf{1,188,074 transactions} and is strongly associated with fraud.

\begin{lstlisting}
zero_after_transfer = df[
    (df["oldbalanceOrg"] > 0)
    & (df["newbalanceOrig"] == 0)
    & (df["type"].isin(["TRANSFER", "CASH_OUT"]))
]
\end{lstlisting}

% ─────────────────────────────────────────────
\section{Feature Engineering}
% ─────────────────────────────────────────────

Two derived features were constructed to capture balance movement:

\begin{align}
\texttt{balanceDiffOrig} &= \texttt{oldbalanceOrg} - \texttt{newbalanceOrig} \\
\texttt{balanceDiffDest} &= \texttt{newbalanceDest} - \texttt{oldbalanceDest}
\end{align}

\begin{itemize}
    \item A \textbf{positive} \texttt{balanceDiffOrig} means money left the sender (expected in fraud).
    \item A \textbf{negative} value (1,399,253 cases) may indicate balance top-ups or data anomalies.
    \item \texttt{balanceDiffDest} had 1,238,864 negative cases, indicating the receiver's balance decreased (suspicious in context of TRANSFER fraud).
\end{itemize}

\subsection{Final Feature Set for Modelling}

\begin{center}
\begin{tabular}{ll}
\toprule
\textbf{Feature} & \textbf{Type} \\
\midrule
\texttt{step}            & Numerical \\
\texttt{type}            & Categorical (One-Hot Encoded) \\
\texttt{amount}          & Numerical \\
\texttt{oldbalanceOrg}   & Numerical \\
\texttt{newbalanceOrig}  & Numerical \\
\texttt{oldbalanceDest}  & Numerical \\
\texttt{newbalanceDest}  & Numerical \\
\texttt{balanceDiffOrig} & Numerical (engineered) \\
\texttt{balanceDiffDest} & Numerical (engineered) \\
\bottomrule
\end{tabular}
\end{center}

Note: \texttt{nameOrig}, \texttt{nameDest}, and \texttt{isFlaggedFraud} are excluded — account IDs are high-cardinality identifiers and \texttt{isFlaggedFraud} would constitute data leakage in a real deployment.

% ─────────────────────────────────────────────
\section{Model Architecture}
% ─────────────────────────────────────────────

\subsection{Pipeline Design}

A scikit-learn \texttt{Pipeline} is used to chain preprocessing and classification:

\begin{lstlisting}
from sklearn.pipeline import Pipeline
from sklearn.compose import ColumnTransformer
from sklearn.preprocessing import StandardScaler, OneHotEncoder
from sklearn.linear_model import LogisticRegression

preprocessor = ColumnTransformer(transformers=[
    ('num', StandardScaler(), numerical_features),
    ('cat', OneHotEncoder(), ['type'])
])

pipeline = Pipeline([
    ('preprocessor', preprocessor),
    ('classifier', LogisticRegression())
])
\end{lstlisting}

\subsection{Preprocessing Steps}

\begin{itemize}
    \item \textbf{StandardScaler} — standardises numerical features to zero mean and unit variance, preventing features with large magnitudes (e.g.\ \texttt{amount}) from dominating the model.
    \item \textbf{OneHotEncoder} — converts the categorical \texttt{type} column into binary indicator columns (one per transaction type).
\end{itemize}

\subsection{Classifier}

\textbf{Logistic Regression} is chosen as the baseline classifier. It models the log-odds of fraud as a linear combination of features:

\begin{align}
P(\texttt{isFraud} = 1 \mid \mathbf{x}) = \frac{1}{1 + e^{-(\mathbf{w}^\top \mathbf{x} + b)}}
\end{align}

\subsection{Train--Test Split}

\begin{lstlisting}
X_train, X_test, y_train, y_test = train_test_split(
    X, y, test_size=0.2, random_state=42, stratify=y
)
\end{lstlisting}

A stratified 80/20 split preserves the fraud ratio in both partitions, which is critical given the 0.13\% class imbalance.

\subsection{Model Persistence}

The trained pipeline is serialised with \texttt{joblib}:

\begin{lstlisting}
import joblib
joblib.dump(pipeline, "fraud_detection_model.pkl")
\end{lstlisting}

% ─────────────────────────────────────────────
\section{Streamlit Web Application}
% ─────────────────────────────────────────────

\subsection{Application Structure}

The app (\texttt{fraud\_detection.py}) provides a browser-based interface for real-time prediction.

\begin{lstlisting}
import streamlit as st
import pandas as pd
import joblib

def load_model(path="fraud_detection_model.pkl"):
    try:
        return joblib.load(path)
    except Exception as e:
        st.error(f"Failed to load model: {e}")
        return None
\end{lstlisting}

\subsection{User Inputs}

The user provides the following fields via the sidebar form:

\begin{center}
\begin{tabular}{ll}
\toprule
\textbf{UI Label} & \textbf{Mapped Column} \\
\midrule
Transaction Type     & \texttt{type} \\
Step (time unit)     & \texttt{step} \\
Amount               & \texttt{amount} \\
Old Balance (Sender) & \texttt{oldbalanceOrg} \\
New Balance (Sender) & \texttt{newbalanceOrig} \\
Old Balance (Receiver) & \texttt{oldbalanceDest} \\
New Balance (Receiver) & \texttt{newbalanceDest} \\
\bottomrule
\end{tabular}
\end{center}

\subsection{Derived Features (Auto-computed)}

The two engineered features are computed automatically from the user inputs before prediction:

\begin{lstlisting}
balanceDiffOrig = oldbalanceOrg - newbalanceOrig
balanceDiffDest = newbalanceDest - oldbalanceDest
\end{lstlisting}

\subsection{Prediction Output}

\begin{itemize}
    \item If \texttt{predict()} returns 1: \textcolor{red}{\textbf{Fraudulent transaction}} (displayed as a red error banner).
    \item If \texttt{predict()} returns 0: \textcolor{green!60!black}{\textbf{Legitimate transaction}} (displayed as a green success banner).
    \item The fraud probability (\texttt{predict\_proba}) is shown alongside the label.
\end{itemize}

\subsection{Running the App}

\begin{lstlisting}[language=bash]
# Activate virtual environment (Windows)
.venv\Scripts\Activate.ps1

# Launch Streamlit
streamlit run fraud_detection.py
\end{lstlisting}

% ─────────────────────────────────────────────
\section{Project File Structure}
% ─────────────────────────────────────────────

\begin{lstlisting}[language=bash, numbers=none]
FraudDetection/
|-- AIML Dataset.csv            # Raw transaction data (6.3M rows)
|-- analysis_model.ipynb        # EDA and model training notebook
|-- fraud_detection_model.pkl   # Serialised trained pipeline
|-- fraud_detection.py          # Streamlit prediction app
|-- fraud_detection_report.tex  # This report
|-- .venv/                      # Python virtual environment
\end{lstlisting}

% ─────────────────────────────────────────────
\section{Dependencies}
% ─────────────────────────────────────────────

\begin{center}
\begin{tabular}{ll}
\toprule
\textbf{Package} & \textbf{Purpose} \\
\midrule
\texttt{pandas}        & Data loading and manipulation \\
\texttt{numpy}         & Numerical operations \\
\texttt{matplotlib}    & Plotting (EDA) \\
\texttt{seaborn}       & Statistical visualisation \\
\texttt{scikit-learn}  & ML pipeline, preprocessing, Logistic Regression \\
\texttt{joblib}        & Model serialisation/deserialisation \\
\texttt{streamlit}     & Interactive web application \\
\bottomrule
\end{tabular}
\end{center}

Install all dependencies into the virtual environment:

\begin{lstlisting}[language=bash, numbers=none]
.venv\Scripts\pip install pandas numpy matplotlib seaborn scikit-learn joblib streamlit
\end{lstlisting}

% ─────────────────────────────────────────────
\section{Key Findings and Observations}
% ─────────────────────────────────────────────

\begin{enumerate}[label=\textbf{\arabic*.}]
    \item \textbf{Extreme class imbalance} (0.13\% fraud) means accuracy alone is a misleading metric. Precision, recall, F1, and AUC-ROC are more informative.

    \item \textbf{Fraud is confined to TRANSFER and CASH\_OUT} transactions. This could be leveraged as a hard filter in a production rule-based pre-screener.

    \item \textbf{Zero-balance-after-transfer} is a strong fraud signal: 1,188,074 transactions match this pattern, many overlapping with known fraudulent cases.

    \item \textbf{Balance difference features} (\texttt{balanceDiffOrig}, \texttt{balanceDiffDest}) explicitly encode the magnitude and direction of money movement, giving the linear model a stronger fraud signal than raw balances alone.

    \item \textbf{Logistic Regression} provides an interpretable baseline. The sign and magnitude of each learned weight directly indicate whether a feature increases or decreases fraud probability.
\end{enumerate}

% ─────────────────────────────────────────────
\section{Possible Improvements}
% ─────────────────────────────────────────────

\begin{itemize}
    \item \textbf{Class imbalance handling} — apply SMOTE oversampling or \texttt{class\_weight='balanced'} in Logistic Regression to improve recall on the minority class.
    \item \textbf{Non-linear models} — Random Forest or XGBoost typically outperform Logistic Regression on tabular fraud data.
    \item \textbf{Threshold tuning} — adjust the decision threshold from 0.5 to maximise F1 or recall depending on business cost of false negatives.
    \item \textbf{Feature selection} — use feature importance or coefficient magnitude to remove low-signal inputs.
    \item \textbf{Cross-validation} — replace the single train--test split with k-fold CV for more robust performance estimates.
\end{itemize}

\end{document}
